\chapter{What's ungrammatical}


\begin{itemize}
    \item There's a common way to say something, a construction that pairs meaning with form.
    \item When something fails to follow this, we assume some new function.
    \item If no such function can be found, it feels ungrammatical.
\end{itemize}

%It's easier to say what is not grammatical than to say what is. Things may be ungrammatical in a relative or absolute sense, relative to a particular language, or simply not possible in human language \citep{Leivada2020} eg. \textit{The doctor the nurse the hospital had hired met John.} (Frazier 1985)


\section{\textit{I go there yesterday}}

\textit{I go there yesterday} sets up a clash between the present tense of the verb, which signals non-past time, and the past-time meaning of the modifier \textit{yesterday}.

\ea\label{ex:go-yesterday-clash}
    \begin{dialogue}
        \speak{Son} \phantom{..*}\textit{Hi, Mom.}
        \speak{Mom} \phantom{*}\textit{Hi, Honey! Where were you?}
        \speak{Son} \phantom{..*}\textit{The new bookstore.}
        \speak{Mom} *\textit{Oh, \myuline{I go there yesterday}.}
    \end{dialogue}
\z

Our default assumption is that Mom isn't just messing up here. We attribute intentionality to her choice of present-tense \mention{go} + \mention{yesterday}. Assuming a clash like this is intended, though, sends Son (and us) on a search for a reason, a quest that hits a dead end, as no justification presents itself. And so we register the underlined clause as ungrammatical.

%And this cost us. Our pupils dilated reflecting the mental effort to resolve the clash. Our heart rate variability dropped.

But could \textit{I go there yesterday} ever be grammatical?

\ea\label{ex:go-yesterday-effective}
    \begin{dialogue}
        \speak{Son} \phantom{..*}\textit{Hi, Mom.}
        \speak{Mom} \phantom{*}\textit{Hi, Honey! Where were you?}
        \speak{Son} \phantom{..*}\textit{The new bookstore.}
        \speak{Mom} \phantom{*}\textit{Oh, you won't believe this. So, \myuline{I go there yesterday}, and\dots}
    \end{dialogue}
\z

In this iteration, Mom offers a vital interpretive key: she is initiating a story. The present-tense form \mention{go}, clashing with the past-oriented \mention{yesterday}, is now reconceptualized as a story-telling technique. We're able to see an -- likely \mention{the} -- intended purpose behind it: a sense of immediacy and involvement in the storytelling, making what was sharply ungrammatical in (\ref{ex:go-yesterday-clash}) indisputably grammatical in (\ref{ex:go-yesterday-effective}).

\begin{center}
    -- --
\end{center}

As Old French became Middle French in the fourteenth century and gradually evolved into Modern French, the present perfect was becoming a past tense in a striking example of various pressures and motivations allowing the ungrammatical to become grammatical.

Until the fourteenth century, French and English had matching simple-past and present-perfect tenses. The simple past looked like this: \textit{Je \uline{changai} mon avis} `I changed my mind', and the present perfect like this: \textit{J'\uline{ai changé} mon avis.} `I have changed my mind'. Though these mean basically the same thing, \mention{changai} `changed' is a past tense while the \mention{ai} of \textit{J'\uline{ai changé}} is the present `have'. That means \textit{J'\uline{ai changé}} is a present tense, albeit a retrospective one. Like the English present perfect, it's anchored in the now but encompasses things past. This makes it very compatible with a word like \mention{recently}, which covers the present and not-too-distant past.

A word like \mention{hier} `yesterday', though, whose meaning is completely separated from the here and now sets up a clash. In Old French, *\textit{J'\uline{ai} changé mon avis \uline{hier}} meant both present (\mention{ai}) and not present (\mention{hier}), making it as ungrammatical as *\textit{I have changed my mind yesterday} in English.

%The `present' meaning of the present perfect can also be seen in discussing event sequencing. Using the past tense, we might say, \textit{He left, and then she arrived,}. But employing the present perfect in similar contexts *\textit{He has left, and then she has arrived} introduces a clash. It suggests his departure is both connected to the present and yet set off from it by her arrival.

To say that the Old-French present perfect turned into a past tense, then, is to say that it lost that `present' meaning, becoming fully compatible with a word like \mention{hier} `yesterday'. In Modern French it's perfectly grammatical to say \textit{J'ai changé mon avis \uline{hier}} because it now means `I changed my mind yesterday,' with no taint of the present. In short, what was the present perfect became the \textsc{passé composé} `the compound past tense'.

Although the meaning of \mention{j'ai changé} has changed, the form really hasn't. It still includes an auxiliary verb, usually \mention{avoir} `have' in its present-tense form, along with the past participle of the verb to be put into the past tense, here \mention{changer} `change'. The full set of forms is shown in Table \ref{tab:have-changed-conjugation}.

\begin{table}[ht]
\centering
\caption{Formation of the Passé Composé in Modern French}
\label{tab:have-changed-conjugation}
\begin{tabular}{llll}
\toprule
\textbf{Subject} & \textbf{\textit{Avoir}} & \textit{\textbf{Changer}} & \textbf{Translation} \\
\midrule
\textit{j'} & \textit{ai} & \multirow{6}{*}{\raisebox{60pt}{\rdelim\}{6}{10pt}[{}]}\raisebox{26pt}{\textit{changé}}} & `I have changed' \\
\textit{tu}      & \textit{as}  &      & `You have changed' \\
\textit{il/elle} & \textit{a} &    & `He/She has changed' \\
\textit{nous}    & \textit{avons} &    & `We have changed' \\
\textit{vous}    & \textit{avez} &     & `You have changed' \\
\textit{ils/elles} & \textit{ont} &    & `They have changed' \\
\bottomrule
\end{tabular}
\end{table}


\bigskip


There's no simple semantic or syntactic explanation for how this set of verb forms changed its meaning. But when social and phonological factors are accounted for, a potential story begins to emerge.

In Old French, as in its ancestor Latin, stress was typically placed on the penultimate syllable. Over centuries, this resulted in an erosion of the final syllable. For example, Latin \textit{fundaMENtum} became Old French \textit{fondeMENT}, preserving only the $\langle$t$\rangle$ of the last syllable. Even that has been lost in Modern French pronunciation /fɔ̃d.mɑ̃/, though the fossilized spelling remains. This kind of \enquote{Matthew effect}\footnote{From the Gospel of Matthew 25:29, \enquote{For to everyone who has, more will be given\dots but from the one who has not, even what he has will be taken away.}} where unstressed elements fade away is pervasive and a common driver of change across languages.

Erosion at the end of words is a problem for Latinate languages because that's where the verb conjugations are, and those conjugations are numerous. The simple past tense's troubles start here. Table \ref{table:past_tense_conjugation} shows the well-marked Old French simple past verb endings and their diminished state in Modern French.

\begin{table}[ht]
\centering
\begin{tabular}{lll}
\hline
\textbf{Person} & \textbf{Old Fr. (IPA)} & \textbf{Modern Fr. (IPA)} \\
\hline
1st Singular & /naviˈʒai/ & /na.vi.ʒe/ \\
2nd Singular  & /naviˈʒas/ & /na.vi.ʒa/ \\
3rd Singular  & /naviˈʒa/ & /na.vi.ʒa/ \\
1st Plural &  /naviˈʒames/ & /na.vi.ʒɑ̃/ \\
2nd Plural &  /naviˈʒastes/ & /na.vi.ʒat/ \\
3rd Plural &  /naviˈʒaʀənt/ & /na.vi.ʒɛʀ/ \\
\hline
\end{tabular}
\caption{Simple Past Tense Conjugation in Latin, Old French, and Modern French}
\label{table:past_tense_conjugation}
\end{table}

These sound shifts were influenced by various factors. One might have been the social life of teen girls, a demographic that often plays a significant role in accelerating and popularizing linguistic change \citep{labov2001}. Their desire to express their unique identities, their experimentation with language, and their influence within tight-knit social circles often make them drivers of linguistic innovation. Imagine a group of girls subtly tweaking the established stress patterns, perhaps for emphasis within their conversations or simply as a playful new way of speaking. These variations may have caught on within their circle, eventually spreading more broadly.

Whatever the vector, stress shifted, and the waning syllables lost their final consonants like so many autumn leaves. Not in writing, but in speech. Because Old French relied on verb endings to distinguish between, for example, first-person singular and third-person plural within a given tense, not to mention distinguishing among tenses, things got complicated. Pronouns, previously optional as in Modern Italian or Spanish, began to gain importance.

For Old French, words were generally pronounced as they were spelled. The Old French pronunciations in the table are somewhat speculative, but they're reasonable approximations for the verb \mention{aimer} `love'. This clearly shows that the Old French conjugations are much more distinct than the Modern ones.

\begin{table}[h]
\centering
\begin{tabular}{lll}
\hline
\textbf{Person}          & \textbf{Old French} & \textbf{Modern French} \\ \hline
1st Singular             & \textit{aim} /aim/   & \textit{j'aime} /ʒɛm/           \\
2nd Singular             & \textit{aimes} /aiməs/ & \textit{tu aimes} /ty ɛm/    \\
3rd Singular             & \textit{aime} /aimə/ & \textit{elle aime} /ɛl ɛm/ \\
1st Plural               & \textit{aimons} /aimɔ̃s/ & \textit{nous aimons} /nu zɛ.mɔ̃/ \\
2nd Plural               & \textit{aimez} /aiməz/ & \textit{vous aimez} /vu zɛ.me/ \\
3rd Plural               & \textit{aiment} /aimənt/ & \textit{ils aiment} /il zɛm/ \\ \hline
\end{tabular}
\caption{Comparative Verb Conjugations in Old French and Modern French with IPA Transcriptions}
\end{table}

As a result, many elements of the simple past began to sound very much alike. Luckily, the irregular nature of the auxiliary verbs for the present perfect (before it became reconceptualized as the passé composé), meant that they were far less subject to this phonemic erosion, and so far less subject to conjugational confusion.

\begin{table}[h]
\centering
\begin{tabular}{lll}
\hline
\textbf{Person} & \textbf{Passé Composé} & \textbf{IPA} \\
\hline
\textit{j'} & \textit{ai aimé} &/e ɛ.me/ \\
\textit{tu} & as aimé &/a ɛ.me/ \\
\textit{il}/\textit{elle} & \textit{a aimé } & /a ɛ.me/ \\
\textit{nous} & \textit{avons aimé} & /a.vɔ̃ ɛ.me/ \\
\textit{vous} & \textit{avez aimé} & /a.ve ɛ.me/ \\
\textit{ils}/\textit{elles} & \textit{ont aimé} & /ɔ̃ ɛ.me/ \\
\hline
\end{tabular}
\caption{Passé Composé Conjugation of \mention{aimer} with IPA Transcriptions}
\end{table}

Since the passé simple and the present perfect could generally be used to refer to the same events, as long as a specific time word like \mention{hier} `yesterday' wasn't used, people started to avoid the more phonologically ambiguous passé simple. If that meant being vague about when something had happened, so be it.

Though this series of changes took roughly 300 years in Paris, they spread much more slowly in the surrounding areas, the result being that the use of the passé simple began to take on a whiff of provincialism, at least in certain uses. Once this had happened, the pressure for Parisians to avoid it was enough to overcome the sense of incompatibility between present nature of the form and the past-nature of the meaning. For them, it had become the passé composé.

\begin{center}
    -- --
\end{center}

We can see just the type of problems that verb-form ambiguity can cause in the famous sentence -- at least relatively famous among linguists -- \textit{The horse raced past the barn fell.} If you've never encountered it before, you likely judged it to be ungrammatical. But I haven't starred it because it only gives the illusion of ungrammaticality, an illusion produced by the overlap between the past tense and the past participle, both of which are \textit{raced}.

Past-tense forms are generally more common than past participles, and this is particularly true directly after a noun; the past participle usually follows an auxiliary such as \textit{be} or \textit{have}, as in \textit{the horse \uline{was} raced}. And so we are tricked into believing that \textit{raced} is the past-tense form. But there's a structure in which the participle doesn't need an auxiliary: a post-head modifier, as in, \textit{the horse \uline{seen} on the hill}. Because \textit{see} distinguishes \textit{saw} from \textit{seen} the illusion collapses. In case it's still not clear, the original sentence was \textit{the horse \op which was\cp~raced past the barn fell}.

Despite this occasional overlap between the past tense and past participle in English, our present perfect hasn't become a past tense. A true past tense, like the passé composé or the simple past in English (again the \enquote{preterite}), means `completed before now'. The English perfect, though, means `looking back from a point', and the \textit{present} perfect establishes that point as `now'. We can often freely choose between the two for any given past event: \textit{I've seen her} or \textit{I saw her}, but *\textit{I've seen her yesterday} is still impossible. The `past-time' meaning of that explicit \textit{yesterday} clashes with the `present-time' meaning of present-tense \textit{have}, and this kind of meaning clash gives rise to feelings of ungrammaticality.
% EDITORIAL-HPC: "Meaning clash" is right only if the clash is inside the
% morphosyntactic form-value coupling. Later distinguish this from lexical
% contradiction, which may produce puzzlement without grammaticality failure.

This leads to the question of why we don't get the same ungrammatical feeling when the paradox is introduced by word choice rather than word forms. Statements such as \mention{the water was so dry} or \mention{it was too light to pick up} provoke curiosity rather than discomfort.

Perhaps, lexical choices like this open up many possible interpretations, such as metaphor, exaggeration, and humour, and so we assume that the paradox was intended. The use of the perfect with a past-time reference like \textit{yesterday}, though doesn't suggest such a purpose, and so immediately registers as a grammatical error.

Consider *\textit{I have two cat.} The intended sentence could be \textit{I have one cat} (wrong word) or \textit{I have two cats} (wrong word form), but my guess is that no editor would change \textit{two} to \textit{one}; they would assume the error is in the grammar. Similarly, \textit{I've seen it yesterday} could be \textit{I've seen it recently} (wrong word) or \textit{I saw it yesterday} (wrong tense). Again, I assume that most editors would blame the tense.

In agreement errors like \textit{a group of students are coming}, we assume that the verb form is incorrect rather than that \textit{group} should be \textit{number} or \textit{majority}. In an example like \textit{give the form from Sarah}, lacking context, we set aside the possibility that the lexical word \mention{give} should have been \mention{get} and assume that the \enquote{grammatical word} \mention{from} should be \mention{to}.

In part, this may be because the degrees of freedom are almost always much smaller with grammar than with lexis. To make this concrete, consider that, across the world's languages, the system of number may mark plurality as it does in English -- singularity being unmarked -- or duality and plurality, as in Arabic. A very few, such as Wamesa, mark triality on pronouns. % TODO: cite a typology source
A few other systems exist, such as marking paucality (the state of being few in number), but this mostly exhausts the possibilities for marking grammatical number, apart from not marking number at all. Now consider tallying the possible nouns to which any system of number applies. The number will be thousands of times greater.

Think of language as a set of Lego blocks. You have a limited number of shapes, and they can only connect in specific ways -- the grammar. If something doesn't fit, you would default to shifting, turning, or trying to connect it in one of those ways rather than assuming you've got the wrong piece. You'd assume you made a placement or \enquote{grammar} mistake rather than a selection or \enquote{lexical} mistake. This reflex is born of an understanding that the grammar rules, much like the physical constraints of Lego, are relatively inflexible. A block misplaced by a few degrees or millimetres simply won't connect. But as long as you've got the angle and position right, you've got a lot of latitude about which blocks you use.

This balance isn't inevitable. In constructing an Ikea desk, the \enquote{vocabulary} of pieces is roughly as constrained as the \enquote{syntax} of combinations. So, if something's wrong, you might as well guess you've got the wrong piece as that you've put it together the wrong way. Languages, though, are much more like Lego sets than Ikea kits, at least in this regard.

The problem of statistical remoteness belongs with impossible languages, where it can do more work. Chapter 10 returns to Chomsky's \mention{colorless green ideas} and Pereira's reply.

Recall \textit{So, I go there yesterday \dots}, from the conversation between the mother and son? I claimed that this use of the present tense for past time didn't feel ungrammatical because the motivation is clear. Conversely, past tense for present time is also possible, as long as it's motivated.

Before going on, we should address the fallacy of monosemy, the mistaken idea that any given word, tense, or structure must have a single \enquote{true} meaning -- \mention{mono-} for single and \mention{semy} from \mention{sēma} `sign or portent'. Quite the contrary! In language, polysemy is the rule.

A familiar example of this fallacy can be found in the self-appointed grammar maven invidiously pretending to misunderstand the meaning of \mention{can} in requests. \textit{Can I leave now} is never a question about ability, and the choice of \textit{may I leave now} isn't about the correct meaning. Rather, \mention{may}, in this question, is a sign that the request is a formal one and an acknowledgement of the greater power of the person being asked. As with the Parisian discomfort with the provincialism of the passé simple, the insistence on \mention{may} for requests is a way to assert the social status of the person being asked, a way to insist on deference.

So, it's a mistake, then, to assume that words are monosemous, and the same is true of tenses. Just as \mention{may} can signify social standing along with permission, possibility, or purpose (\textit{You \uline{may} go}; \textit{It \uline{may} rain}; \textit{\dots~that they \uline{may} prosper}), the past tense can convey deference, formality, likelihood, and impossibility (\textit{\uline{Could} you help me?}; \textit{\uline{Did} you wish to reconsider?}; \textit{If I \uline{did} go, \dots}; \textit{If she \uline{was} a client, \dots}). Perhaps the common element motivating these seemingly diverse meanings is the idea of distance, temporal, social, possible, or factual.

Compare (\ref{ex:am-was-wondering}a) with (\ref{ex:am-was-wondering}b).

\ea \label{ex:am-was-wondering}
    \ea[]{\textit{I'm wondering if this is a good idea.}\hfill[\textsc{present tense}]}
    \ex[]{\textit{I was wondering if this was a good idea.}\hfill[\textsc{past tense}]}
    \z
\z

It's difficult to set aside our learned associations between the past tense and tentativeness, but looking as strictly as we can at the basic time meaning, (a) implies a current discomfort with the idea, while (b) suggests that discomfort was in the past and may have been overcome. As a result, (b) is less \enquote{pushy} than (a), more open to allowing the listener their own opinion. This kind of example then provides a motivation for interpreting the past tense more generally as less direct, as more polite.

Another motivation for using the past tense to indicate politeness is that it provides a kind of point-of-view shift. By using the past tense, the speaker \enquote{moves \textit{as if} into the future} \citep[204]{Brown1978}, thus making the statement or question less immediate, less of an imposition on the listener.

These uses are so appealing that, in some cases, the past-tense forms of modal auxiliary verbs have left behind their past-time meanings. \mention{Could}, as the past-tense of \mention{can} `is able to', rarely means `was able to' in today's English, though that was its typical meaning for hundreds of years. The use of \mention{could} for requests, which seems so natural to us, appears to be an innovation of Modern English. % TODO: verify against OED
According to the \textit{Oxford English Dictionary}, the first recorded example of \mention{could} used in a request is in a letter from Lord Chesterfield in 1748, \textit{Could you send me\dots~some seed of the right Cantelupe melons?} Once the request use got going, though, it took off and took over. Presumably, its politeness meaning was so obvious that it wasn't felt to be ungrammatical, at least, not by many.

\begin{center}
    -- --
\end{center}

To reiterate, grammar conveys meaning, and when the meaning clashes with the meaning of some other element of the sentence, we experience that clash as ungrammaticality. We don't get that same feeling when the meanings of words appear to clash mostly because we're used to words telling us stuff we're not used it. At the same time, differences in phonology don't provoke feelings of ungrammaticality because phonology is extremely variable and is generally not a carrier of semantic meaning.

Here's another example of a form-meaning mismatch that results in ungrammaticality, one that, as far as I know, has never been described, and one that seems difficult to attribute solely to structural factors.

A few days ago, as I was working on this chapter, my friend Geoff Pullum emailed me to comment on the strange sentence in (\ref{ex:Rubyar-bad}) that had been published online by the BBC. The star is Geoff's, but I completely agree that this is ungrammatical.

\ea[*]{\textit{The Rubymar is the ship to have been sunk by the Houthis.}}\label{ex:Rubyar-bad}
\z

What, he wondered, could be wrong with this sentence. The hypothesis he suggested was that the \mention{to have been \dots} part needs something to \enquote{license} it. What he had in mind was a process like the one that allows the \mention{than} phrase in cases like \textit{It is more radical a plan \uline{than that}.} Here, \mention{more} licenses \mention{than that}. Without \mention{more}, you get the completely awful *\textit{It's radical a plan than that}. And Geoff wondered if \mention{to have been \dots} might need a word like \mention{first}, \mention{only}, or \mention{largest} to render the sentence grammatical as in (\ref{ex:Rubyar-fixed}), just as \mention{than \dots} needs \mention{more}.

\ea \label{ex:Rubyar-fixed}
    \ea[]{\textit{The Rubymar is the \uline{first} ship to have been sunk by the Houthis.}}
    \ex[]{\textit{The Rubymar is the \uline{only} ship to have been sunk by the Houthis.}}
    \ex[]{\textit{The Rubymar is the \uline{largest} ship to have been sunk by the Houthis.}}
    \z
\z

One problem this poses is that these three words aren't all the same grammatically. Sure, they're all adjectives, but consider that neither \mention{large} nor \mention{larger} work, and they're just as adjective as \mention{largest}. This suggests that it might be the superlative suffix \mention{-est} that's doing the licensing. With this hypothesis in mind, we can experiment with other superlatives to see how they fare. And, indeed, they fare very well: \mention{the best ship}, \mention{the fastest}, \mention{the strongest}, \mention{the most expensive}, \mention{the most recent}, etc.

But that doesn't explain why \mention{first} and \mention{only} work. While only \mention{largest} is a true superlative, the other two words only hint at superlative qualities. \mention{First} can imply the `most important' or the `most desirable' position, and \mention{only} hints at being `the best' or `the most exclusive'. And when it comes to the forms, \mention{first} is already the base form, like \mention{large}, not like the superlative form \mention{largest}. \mention{First} isn't the superlative of a base *\mention{firs} and comparative *\mention{firser}. Nor is \mention{only} the superlative of \mention{on} or anything else. There's no word from which it could be inflected, no word indicating a state of \mention{limited} exclusivity or a \mention{preliminary} form of uniqueness, ideas that seem entirely paradoxical. So, it can't be that the superlative is licensing \mention{to have been \dots} in the way that the comparative licenses the \mention{than} phrase.

Another problem for the licensing hypothesis is that we can construct similar sentences that work without any apparent licensing words. Before moving on, let's look at the structure to see what Geoff and I were considering \enquote{similar sentences}.

The key part seemed to be the noun phrase \textit{the ship to have been \dots}. So we have the definite article \textit{the} plus a head noun, and then a modifier in the form of a \textit{to} infinitival. This includes \textit{to} plus \textit{have} to mark the perfect, along with a past participle. In this case, the past participle was \textit{been}, but it could just as well have been \textit{gone}, \textit{seen}, \textit{found}, \textit{made}, \textit{called}, \textit{smoked}, \textit{minced}, or \textit{barbecued}. So, the part we're interested in is \uline{\textit{the} \textsc{noun} \textit{to have} \textsc{past participle}}.

And so, I went looking for sentences of this kind and found a small number.\footnote{You can see the results at https://www.english-corpora.org/coca/?c=coca\&q=119073728.} This was surprising, because we'd felt that there was something wrong with the structure. And so we decided to zoom out a bit to see what context the examples appeared in. But when I prefaced the search with \mention{is}, the tool I was using wouldn't allow a search that long with so many common words.

To get around this, I went back to the previous list of results and started browsing the full concordance of sentences. As is common when you do this kind of research, you find that lots of matching strings aren't also matching phrases. \textit{Is it a mistake \ob for \uline{the president\cb to have joined} the lawsuit} has the \textit{to} infinitive as the complement of \textit{mistake}, not the modifier of \textit{president}. \textit{City records show the officer alleged \ob in \uline{the lawsuit\cb to have fired} the fatal shots joined the CPD in 2014} has it as a complement of \textit{alleged}, etc.

Eventually, though, I turned up \textit{She seems to be the one to have changed the most,} which seemed fine to both of us. At this point, the whole premise that there was something structurally wrong with \uline{\textit{the} \textsc{noun} \textit{to have} \textsc{past participle}} was in tatters. After staring at it for some time, I realized something new: \textit{is} is present tense, while \textit{be} has no tense at all. It's infinitival. But, we asked ourselves, why should (\ref{ex:seems-to-be-vs-is}a) be better than (\ref{ex:seems-to-be-vs-is}b).

\ea \label{ex:seems-to-be-vs-is}
    \ea[]{\textit{She seems to be the one to have changed the most.}}
    \ex[*]{\textit{She is the one to have changed the most.}}
    \z
\z

\noindent I also found examples (\ref{ex:start}--\ref{ex:stop}).

\ea[]{\textit{Again, he \uline{was} the one to have handled it correctly.}}\label{ex:start}\z
\ea[]{\textit{\dots you \uline{were} probably also the type to have memorized them}}\z
\ea[]{\textit{He should have \uline{been} the one to have run for president}}\z
\ea[]{\textit{Petraeus is said to have \uline{been} the one to have broken off the extramarital affair}}\z
\ea[]{\textit{Mrs. Stanton \uline{was} not the type to have \enquote{stayed up all night dancing}}}\z
\ea[]{[he] \textit{should have \uline{been} the one to have given the land to the boy}}\label{ex:stop}\z

In no case was the (underlined) verb before this phrase in the present tense. This directed our attention back to the \textit{to} infinitival, and we realized that \textit{have} \textsc{past participle} was an infinitival perfect.

The present perfect has already come up a lot in this chapter. As a present tense, it's anchored in the now, but as a perfect, it's retrospective. The infinitive perfect, on the other hand, has no anchor in the present. It's just retrospective, the semantic equivalent of a past tense. And I think the feeling of ungrammaticality is coming from a meaning clash between the present-tense of \textit{is} in (\ref{ex:Rubyar-bad}) and this effective past tense.

This suggests that the problem is a clash between the present tense in the matrix clause and the infinitival perfect in the subordinate. But that doesn't explain why adding in a modifier like \mention{first}, \mention{only}, or \mention{largest} should fix things. (And, yes, I know that \enquote{like} is still undefined here.)

Other things need considering. Interrogatives seem better than declaratives.

\ea
    \ea[?]{\textit{Who is the one to have betrayed the people of India?}}
    \ex[?]{\textit{Is he the one to have betrayed the people of India?}}
    \ex[*]{\textit{He is the one to have betrayed the people of India.}}
    \ex[]{\textit{He is the one who betrayed the people of India.}}
    \z
\z

Also if you change the token to a type, you can coerce it into grammaticality:

\ea
    \ea[*]{\textit{She's the one to have read a handful of books from a century ago.}}
    \ex[]{\textit{She's the type to have read a handful of books from a century ago.}}
    \z
\z

\section{The Ex-Lax conundrum}

% TODO: rewrite — AI-tic transition
Sometimes, what appears ungrammatical at first glance can become perfectly acceptable given the right context. To illustrate this point, let's consider a striking example from the Nixon years in the US. This example, which appears in a dialogue between James McCawley and Herman  \citet[252]{Parret1974}, builds on insights about the grammaticality of responses to questions, an idea developed by Jerry \citet{Morgan1973}. Together, these contributions illuminate the complex relationship between grammar and context, foreshadowing the kinds of issues we'll encounter in the next section.

McCawley presented the following sentence:

\ea[*]{\textit{Spiro conjectures Ex-Lax.}}\label{ex:spiro-ex-lax}\z

At first glance, most syntacticians would agree with the asterisk, marking this as ungrammatical. The verb \textit{conjecture} typically doesn't take a direct object denoting a physical substance. However, Morgan then posed a question that dramatically shifts our perception:

\ea{\textit{Does anyone know what Mrs. Nixon frosts her cakes with?}}\label{ex:pat-nixon-cakes}\z

In the context of this question, the previously ungrammatical sentence suddenly becomes acceptable. We interpret it as an elliptical response, implying a structure more like:

\ea{\textit{Spiro conjectures }[\textit{that Pat Nixon frosts her cakes with}]\textit{ Ex-Lax.}}\label{ex:spiro-full}\z

(In case you're unfamiliar with the historical context, Ex-Lax is a chocolate-flavored laxative, and Spiro Agnew was the Vice President under Richard Nixon. I'll say no more.)

Geoff Pullum reminded me of this example and how it challenged a fundamental assumption in syntax at the time: that it was possible to categorize sentences as grammatical or ungrammatical without considering their communicative context. Morgan's insight was that grammaticality judgments can't be made in isolation from the sentence's intended use and situational context.
% EDITORIAL-HPC: Avoid overgeneralizing Ex-Lax into "context fixes
% grammaticality." Context may reveal a different construction, ellipsis, or
% form-value pairing. It doesn't make every structural violation grammatical.

This example demonstrates how context can dramatically alter our perception of grammaticality. As we'll see in the next section, this principle extends beyond simple question-answer pairs to more complex grammatical structures. The case of the supposedly missing independent relative \mention{whose} provides another striking illustration of how our judgments of grammaticality can be influenced by contextual factors.

\section{What this leaves open}

The case of independent relative \mention{whose} has been moved to Chapter \ref{ch:getting-grammaticality-wrong}, where it can serve as the book's worked example of expert grammaticality judgement going wrong. Here the point is narrower: a sentence may feel ungrammatical because the form and the communicative situation have no obvious way to fit together.
